\chapter{Anal Pain}

\section*{Definition}
Anal pain is pain or discomfort in the anal canal or surrounding skin. It may be sharp, burning, throbbing, or pressure-like and may worsen with bowel movements or sitting. Because several conditions can cause it, pain alone cannot establish the diagnosis.

\section*{When to See a Doctor}

\begin{itemize}
 \item Severe or worsening pain, fever, chills, pus or drainage, a painful lump, heavy bleeding, black stool, or inability to pass stool or urine
 \item Pain after anal trauma or with a new mass, persistent bleeding, weight loss, or a change in bowel habits
\end{itemize}

\section*{How It Develops}
The anal canal and surrounding skin contain many sensory nerves. Acute pain commonly results from a fissure, thrombosed external hemorrhoid, infection, or trauma. Persistent or recurrent pain may reflect fissure, fistula, inflammatory disease, pelvic-floor pain, nerve pain, or a less common tumor. A painful abscess can worsen quickly and needs prompt treatment.

\section*{Causes}
\begin{itemize}
 \item Anal fissure (acute or chronic)
 \item Thrombosed external hemorrhoid
 \item Perianal abscess
 \item Anal fistula
 \item Proctalgia fugax (brief sudden rectal spasm)
 \item Levator ani syndrome (chronic pelvic floor spasm)
 \item Inflammatory bowel disease (Crohn disease, UC)
 \item Proctitis (infectious or radiation-induced)
 \item Pruritus ani with skin breakdown or dermatitis
 \item Anal malignancy (rare)
\end{itemize}

\section*{Related Symptoms}
\begin{itemize}
 \item Rectal bleeding
 \item Constipation
 \item Itching
 \item Swelling
 \item Fever
\end{itemize}

\section*{Medical Evaluation}
\begin{itemize}
 \item The clinician will inspect the perianal area and, when tolerable and appropriate, perform a gentle digital rectal examination. Severe pain may require analgesia, local anesthetic, or specialist assessment before an internal examination.
 \item Anoscopy or proctoscopy may assess internal hemorrhoids, fistulas, bleeding, or proctitis when indicated; it is not necessary for every episode of acute anal pain.
 \item Laboratory studies including complete blood count and inflammatory markers (CRP, ESR) may be ordered if inflammatory bowel disease or infection is suspected.
 \item Referral to a colorectal surgeon is indicated for suspected fistula, abscess requiring drainage, or persistent symptoms despite conservative management.
\end{itemize}

\section*{Treatment Options}
\begin{itemize}
 \item Topical nitroglycerin or calcium channel blockers (nifedipine, diltiazem) relax the internal anal sphincter and promote healing of anal fissures.
 \item Rubber band ligation, sclerotherapy, or infrared coagulation is used for symptomatic internal haemorrhoids.
 \item Perianal abscesses generally require prompt incision and drainage. Antibiotics are added for selected patients, such as those with systemic infection, spreading cellulitis, immune compromise, or certain associated conditions; they do not replace drainage.
 \item Surgical options include fistulotomy for anal fistulas, lateral internal sphincterotomy for chronic fissures, and haemorrhoidectomy for advanced haemorrhoidal disease.
\end{itemize}

\section*{Self-Care and Home Management}
\begin{itemize}
 \item Take warm shallow baths for 10--15 minutes as needed for comfort, and gently pat the area dry.
 \item Increase dietary fiber gradually and drink enough fluid to keep stools soft, unless fluid restriction has been prescribed.
 \item Ask a pharmacist or clinician before using topical hemorrhoid products or anesthetics; use them briefly and stop if irritation develops.
 \item Use a clinician-recommended fiber supplement or laxative if constipation is contributing, and avoid straining or prolonged toilet sitting.
 \item Avoid prolonged sitting on hard surfaces; use a soft cushion if helpful, but avoid devices that increase pressure or worsen pain.
\end{itemize}

\section*{References}
\begin{thebibliography}{99}
\bibitem{anal_pain:ref1} Beaty JS, Shashidharan M. ``Anal fissure.'' \textit{Clin Colon Rectal Surg}. 2016;29(1):30-37.
\bibitem{anal_pain:ref2} Lohsiriwat V. ``Hemorrhoids: from basic pathophysiology to clinical management.'' \textit{World J Gastroenterol}. 2012;18(17):2009-2017.
\bibitem{anal_pain:ref3} Wald A. ``Diagnosis and management of anorectal pain.'' \textit{Am Fam Physician}. 2017;95(4):234-240.
\bibitem{anal_pain:ref4} Beck DE, Roberts PL, et al. \textit{The ASCRS Textbook of Colon and Rectal Surgery}, 3rd ed. Springer, 2018.
\bibitem{anal_pain:ref5} UpToDate. ``Anal fissure, hemorrhoids, and anorectal abscess.'' Wolters Kluwer, 2023.
\bibitem{anal_pain:ref6} National Health Service. Anal pain and anal fissure patient guidance. Updated 2025.
\end{thebibliography}
